% s1-result.tex -- Section 1 of the reader's companion (drafter: aqft, R3; edited in R4;
% revised 7 November 2026 for the record of that date: theorem box, status, ledger, pages).
% Mathematics from part_i/main.tex only; numbers and pages from companion/council/labels.md.
\section{What is proved, and under what}\label{sec:result}

\emph{The full net converges, with a rate; the one GJS theorem imported, Theorem~1.1.7, is
also used in a case its printed statement covers, beyond its printed proof.}

\emph{Setting.}
$\HH$ is $Q$-space over the time-zero free field $\varphi_0$ of mass $m_0>0$ and covariance
$C=(2\mu)^{-1}$, $\mu=(-\partial_x^2+m_0^2)^{1/2}$. For a bounded
open interval $O$, $\A(O)$ is the von Neumann algebra generated by the Weyl operators
$W(f,0)=e^{i\varphi_0(f)}$ and $W(0,h)=e^{i\pi_0(h)}$, $f,h\in C_c^\infty(O)$ real, with
$\pi_0$ the time-zero momentum field, and
$\A$ is the norm closure of $\bigcup_O\A(O)$
\rec{\S2.1}{7}. The interaction sits in the states and in the Glimm--Jaffe dynamics
$\alpha$, a $g$-independent one-parameter automorphism group of $\A$ \rec{A6(ii)}{13},
not in the local algebras. With
$V(h)=\lambda\int h(x)\,{:}\mathscr P(\varphi_0(x)){:}_C\,dx$ (coupling $\lambda$,
$C$-Wick order), where $\mathscr P(\xi)=\sum_{n=0}^{2p'}a_n\xi^n$ is the interaction
polynomial, set
$K(g)=H_0+V(g)$ ($H_0$ free), $E(g)=\inf\spec K(g)$ and $H(g)=K(g)-E(g)$;
$\Om_g$ is its normalised ground vector, $\omega_g=\ip{\Om_g}{\cdot\,\Om_g}$ the cutoff
state, and $\beta^g_t=\mathrm{Ad}\,e^{itH(g)}$. In \rec{Definition~2.1}{8}, $\CE$
and $\CH$ are the measurable $[0,1]$-valued functions of compact support on $\R^2$ and on
$\R$, and $\G\subset\CH$ is the plateau class of $g\in C_c^\infty(\R)$ with $0\le g\le1$
and $g\equiv1$ on a nonempty open interval, directed by inclusion of the plateau interiors
$\{g=1\}^\circ$; $O_r$ is the open $r$-neighbourhood of $O$, and
$d(g,O)=\dist(O,\R\setminus\{g=1\}^\circ)$. Throughout,
$0\le\lambda/m_0^2<\varepsilon_{\mathscr P}:=\lambda_*/M_*^2$, with $M_*,\lambda_*$
chosen so that the hypotheses printed for Theorem~1.1.7 of Glimm--Jaffe--Spencer (GJS)
\cite{GJS} hold at bare mass $M_*$ \rec{\S2.6}{20}.

\emph{The theorem.}
Below, $\tau_a$ are the space translations and $\theta$ is the parity automorphism
\rec{\S2.3}{9}; $m_1>0$ is the GJS decay constant \rec{A8(iv)}{15}; $M$ is the moment
constant of \rec{(M$_2$)}{24}; and $\Psi_{m_1}$ is the rate function $\Psi_\gamma$ of
\rec{(29)}{36} at the gap parameter $\gamma=m_1$, a function of the distance $d$ stated in
\S\ref{sec:lppl}.

\begin{recordbox}{Theorem 12.2 of the record (pp.\,38--39)}
Let $0\le\lambda/m_0^2<\varepsilon_{\mathscr P}$. Then there is a state $\omega_\infty$ on
$\A$ such that:
\begin{enumerate}[leftmargin=2.2em]
\item[\textup{(1)}] \emph{(full-net convergence with rate)} for every bounded open interval
      $O$, $\lim_{g\in\G}\|(\omega_g-\omega_\infty)|_{\A(O)}\|=0$, and
      $\|(\omega_g-\omega_\infty)|_{\A(O)}\|\le\Psi_{m_1}(d)$ whenever $g\equiv1$ on $O_d$,
      with $\Psi_{m_1}$ decaying faster than every inverse power;
\item[\textup{(2)}] \emph{(local normality)} $\omega_\infty|_{\A(O)}$ is normal for every $O$;
\item[\textup{(3)}] \emph{(ground state)} $\omega_\infty$ is a ground state of $(\A,\alpha)$
      in the spectral sense;
\item[\textup{(4)}] \emph{(invariances)} $\omega_\infty\circ\tau_a=\omega_\infty$ for all
      $a\in\R$; and for even $\mathscr P$, $\omega_\infty\circ\theta=\omega_\infty$;
\item[\textup{(5)}] \emph{(spectra and local moment)}
      \[
       \spec H(g)\subseteq\{0\}\cup[m_1,\infty)\quad(g\in\CH),
       \qquad \dim\ker H(g)=1,
      \]
      and $\sup_{g\in\CH}\|V(h)\Om_g\|\le M$ for real $h$ with $\|h\|_\infty\le1$ supported
      in an interval of length $2$.
\end{enumerate}
[\,\dots] The constant $M$ of item (5), and through it the rate $\Psi_{m_1}$ of item (1),
contain the finite, unevaluated constant $S_2(B_{3\times5})$ of Theorem~C.5.
\end{recordbox}

The elided sentences name the one GJS theorem imported, Theorem~1.1.7 with its transverse
case A8(iv)-T, and the inputs of Theorems~D.3--D.4, retire A8(vi), and list A1--A7 with
\rec{Proposition~3.2}{21}. Items (1)--(4) are statements along $\G$; item (5) concerns the
cutoff Hamiltonians $H(g)$, $g\in\CH$.

\emph{The status.}
Theorem~12.2 is proved from printed statements only, with one GJS theorem imported:
Theorem~1.1.7, which is also used for sharp-time observables across strips transverse to
the time-zero line, a case its printed statement covers and its printed proof does not
reach (listed separately as A8(iv)-T). The uniform local moment is a theorem of the record
(Appendix~D). The human readings of audit item C5 that this status rests on, items 1--4
and 10--12 with item 9(c), are open \recpp{title-page notice}{1--2}; see
\S\ref{sec:import}.

\emph{What is new.}
Convergence, not the topology. The unaveraged full net $(\omega_g)_{g\in\G}$ has one limit
on every local algebra $\A(O)$, with a rate, and every $w^*$ limit point of the
construction in Glimm and Jaffe's paper III \cite{GJ3} (GJ3 below) coincides with
$\omega_\infty$ on each $\A(O)$ \rec{Corollary~12.3}{39}. This makes $\omega_\infty$
canonical for the cutoff family. Uniqueness among the locally normal ground states of
$(\A,\alpha)$, in the spectral sense of \rec{Proposition~12.1}{37}, is not claimed
\rec{\S1.3}{5}.

\begin{context}
% (R5)
In GJ3, the averaged states of \cite[(2.4)]{GJ3} lie, on each local algebra, in a
norm-compact subset of the local dual, and their limit points are normal
\cite[Thm.~2.3, p.~211]{GJ3}; in the representation of any $w^*$ limit point, the
generator of time translations has spectrum in $[0,\infty)$ \cite[Thm.~2.1, p.~211]{GJ3}.
\end{context}

Table~\ref{tab:ledger} sorts the statements this companion explains by what their proofs
use; Fig.~\ref{fig:spine} shows where A8(iv)-T enters. Status words have one meaning each:
proved, partial, conjectured, numerical, refuted; \emph{proved} means proved from the
imports A1--A8, read with the open human readings of C5 beside it \recpp{Appendix~B}{44--45}.

\begin{table}[ht]
\centering\small
\begin{tabular}{@{}>{\raggedright\arraybackslash}p{0.24\textwidth}
                >{\raggedright\arraybackslash}p{0.72\textwidth}@{}}
\toprule
status & statements of the record (printed page)\\
\midrule
proved; \Ttext{rests on A8(iv)-T}, directly or through $M$
 & Theorems D.4 (95), 5.1 (24), 11.1 (36), 12.2 (38), 13.2 (40); Corollaries 5.2 (25),
   12.3 (39); Lemmas 10.2 (35), 10.3 (35), 10.5 (36); Propositions 10.4 (35),
   11.2 (37).\\[2pt]
proved; Theorem~1.1.7 across parallel strips only
 & Theorem 7.8 (30), with Lemmas 7.1--7.4 (26--28), Proposition 7.5 (28), Lemma 7.6 (29),
   Theorem 7.7 (29) and Theorem D.3 (90).\\[2pt]
proved, independent of A8
 & Propositions 2.7 (19), 8.1 (31); Lemmas 2.8 (20), 4.4 (23), 9.1 (32), 9.2 (33),
   C.2--C.4 (48--49), D.1 (79); Theorem C.5 (51).\\[2pt]
proved as implications under a stated hypothesis
 & Proposition 8.2 (32) and Lemmas 9.3 (34), 10.1 (34): under the gap hypothesis
   (G$_\gamma$) (30), which for $P(\varphi)_2$ only Theorem 7.8 supplies.
   Proposition 12.1 (37): for any $(g_n)\subseteq\G$ with $d(g_n,O)\to\infty$ for every
   bounded interval $O$ and $\omega_{g_n}$ converging in norm on every $\A(O)$.\\[2pt]
not claimed; open
 & Not claimed: uniqueness of $\omega_\infty$ among locally normal ground states of
   $(\A,\alpha)$ in the spectral sense, a sharp rate, anything outside the weak-coupling
   region (Theorem 13.2 is conditional on Hypothesis 13.1), a gap for a GNS Hamiltonian of
   $\omega_\infty$, identification with the Euclidean (Wightman) theory of GJS, Lorentz
   covariance (not discussed; \S\ref{sec:boundaries}). Not done: a proof of A8(iv)-T, a
   value of the constant of Theorem C.5 (\S\ref{sec:import}). \emph{Open:} the human readings
   of C5 items 1--4 and 10--12, with item 9(c) (\S\ref{sec:record}); the Part~II gate
   (II.11) (\S\ref{sec:part2}).\\
\bottomrule
\end{tabular}
\caption{Claims ledger, read from the record's proofs. Appendix~A classifies every
numbered statement.}
\label{tab:ledger}
\end{table}

\noindent\emph{In the record:} title-page notice, pp.~1--2; \S1.3, p.~5; \S2.1, p.~7;
Definition~2.1, p.~8; \S2.3, p.~9; A6, p.~13; A8(iv), p.~15; \S2.6, p.~20;
Propositions~11.2 and 12.1, Theorem~12.2 and Corollary~12.3, pp.~37--39; Appendix~B,
pp.~44--45.
